\documentclass[11pt,a4paper]{article}

% shared/macros.tex
% Shared notation for all surface-partition math documents.
% Include with: % shared/macros.tex
% Shared notation for all surface-partition math documents.
% Include with: % shared/macros.tex
% Shared notation for all surface-partition math documents.
% Include with: \input{../shared/macros}

% -------------------------------------------------------------------------
% Packages (include here so each document only needs \input{../shared/macros})
% -------------------------------------------------------------------------
\usepackage{amsmath,amssymb,amsthm}
\usepackage{bm}           % \bm for bold math
\usepackage{mathtools}    % \coloneqq, dcases, etc.
\usepackage{booktabs}     % \toprule etc. in tables
\usepackage{hyperref}
\usepackage{geometry}
\usepackage{microtype}
\usepackage{parskip}
\usepackage{xcolor}
\usepackage{tcolorbox}
\usepackage{enumitem}
\usepackage{float}             % [H] float placement

\geometry{a4paper, margin=2.5cm}

% -------------------------------------------------------------------------
% Theorem-like environments
% -------------------------------------------------------------------------
\theoremstyle{definition}
\newtheorem{defn}{Definition}[section]
\newtheorem{remark}{Remark}[section]

\theoremstyle{plain}
\newtheorem{prop}{Proposition}[section]

% -------------------------------------------------------------------------
% Mesh and geometry
% -------------------------------------------------------------------------
\newcommand{\V}{\mathbf{V}}            % vertex array V[i] in R^dim
\newcommand{\F}{\mathbf{F}}            % face (triangle) array
\newcommand{\vone}[1]{v_{1,#1}}       % first vertex index of VP i
\newcommand{\vtwo}[1]{v_{2,#1}}       % second vertex index of VP i

% -------------------------------------------------------------------------
% Variable points
% -------------------------------------------------------------------------
\newcommand{\lam}{\lambda}             % scalar VP parameter
\newcommand{\Lam}{\boldsymbol{\lambda}}% full parameter vector
\newcommand{\vp}[1]{\mathbf{p}_{#1}}  % 3D position of VP i
\newcommand{\ed}[1]{\mathbf{d}_{#1}}  % edge direction d_i = V[v1_i] - V[v2_i]

% -------------------------------------------------------------------------
% Steiner / triple-point
% -------------------------------------------------------------------------
\newcommand{\Sv}{\mathbf{S}}           % Steiner (Fermat-Torricelli) point
\newcommand{\Ki}{K_{i}}               % projection matrix (I-n_i n_i^T)/r_i
\newcommand{\Mmat}{M}                 % sum of K_i matrices at a triple point

% -------------------------------------------------------------------------
% Normals and unit vectors
% -------------------------------------------------------------------------
\newcommand{\nhat}{\hat{\mathbf{n}}}  % unit normal to a triangle
\newcommand{\Dhat}{\hat{\boldsymbol{\Delta}}}  % unit segment direction

% -------------------------------------------------------------------------
% Operations
% -------------------------------------------------------------------------
\newcommand{\norm}[1]{\left\|#1\right\|}
\newcommand{\abs}[1]{\left|#1\right|}
\newcommand{\ip}[2]{\langle #1,\, #2 \rangle}

% Projection perpendicular to a unit vector \uhat
\newcommand{\Proj}[1]{\bigl(\mathbf{I} - #1 #1^\top\bigr)}

% argmax operator (winner-take-all assignment in Phase 1 documents)
\DeclareMathOperator*{\argmax}{arg\,max}
\DeclareMathOperator*{\argmin}{arg\,min}

% -------------------------------------------------------------------------
% Calculus shorthands
% -------------------------------------------------------------------------
\newcommand{\pd}[2]{\dfrac{\partial #1}{\partial #2}}
\newcommand{\pdd}[3]{\dfrac{\partial^2 #1}{\partial #2\,\partial #3}}
\newcommand{\grad}{\nabla}
\newcommand{\Hess}[1]{\nabla^2 #1}

% -------------------------------------------------------------------------
% Optimization problem objects
% -------------------------------------------------------------------------
\newcommand{\Preg}{P_{\mathrm{reg}}}  % regular perimeter
\newcommand{\Pst}{P_{\mathrm{st}}}   % Steiner perimeter
\newcommand{\Areg}{A^{\mathrm{reg}}} % regular cell area
\newcommand{\Ast}{A^{\mathrm{st}}}   % Steiner area contribution
\newcommand{\Atgt}{\bar{A}}           % target area per cell
\newcommand{\Lagr}{\mathcal{L}}       % Lagrangian
\newcommand{\objfac}{\sigma}          % IPOPT objective scaling factor

% -------------------------------------------------------------------------
% Code cross-references (for "In the code..." callouts)
% -------------------------------------------------------------------------
\newcommand{\code}[1]{\texttt{#1}}
\newcommand{\codefile}[1]{\texttt{\small #1}}

% -------------------------------------------------------------------------
% Threshold dynamics / approach B (10-mbo-auction-dynamics)
% -------------------------------------------------------------------------
\newcommand{\Dlump}{D}                 % lumped mass matrix diag(v)
\newcommand{\Atau}{A_\tau}             % discrete diffusion operator (M + tau K)^{-1} M
\newcommand{\Etau}{E_\tau}             % heat-content (threshold-dynamics) energy
\newcommand{\Rcell}{R_{\mathrm{cell}}} % radius of the equal-area disc sqrt(Abar/pi)
\newcommand{\hmean}{h_{\mathrm{mean}}} % mean edge length
\newcommand{\hmax}{h_{\mathrm{max}}}   % max edge length
\newcommand{\Gt}{G_t}                  % heat kernel at time t
\newcommand{\Ktau}{K^{\mathrm{res}}_\tau} % resolvent (backward-Euler) kernel


% -------------------------------------------------------------------------
% Packages (include here so each document only needs % shared/macros.tex
% Shared notation for all surface-partition math documents.
% Include with: \input{../shared/macros}

% -------------------------------------------------------------------------
% Packages (include here so each document only needs \input{../shared/macros})
% -------------------------------------------------------------------------
\usepackage{amsmath,amssymb,amsthm}
\usepackage{bm}           % \bm for bold math
\usepackage{mathtools}    % \coloneqq, dcases, etc.
\usepackage{booktabs}     % \toprule etc. in tables
\usepackage{hyperref}
\usepackage{geometry}
\usepackage{microtype}
\usepackage{parskip}
\usepackage{xcolor}
\usepackage{tcolorbox}
\usepackage{enumitem}
\usepackage{float}             % [H] float placement

\geometry{a4paper, margin=2.5cm}

% -------------------------------------------------------------------------
% Theorem-like environments
% -------------------------------------------------------------------------
\theoremstyle{definition}
\newtheorem{defn}{Definition}[section]
\newtheorem{remark}{Remark}[section]

\theoremstyle{plain}
\newtheorem{prop}{Proposition}[section]

% -------------------------------------------------------------------------
% Mesh and geometry
% -------------------------------------------------------------------------
\newcommand{\V}{\mathbf{V}}            % vertex array V[i] in R^dim
\newcommand{\F}{\mathbf{F}}            % face (triangle) array
\newcommand{\vone}[1]{v_{1,#1}}       % first vertex index of VP i
\newcommand{\vtwo}[1]{v_{2,#1}}       % second vertex index of VP i

% -------------------------------------------------------------------------
% Variable points
% -------------------------------------------------------------------------
\newcommand{\lam}{\lambda}             % scalar VP parameter
\newcommand{\Lam}{\boldsymbol{\lambda}}% full parameter vector
\newcommand{\vp}[1]{\mathbf{p}_{#1}}  % 3D position of VP i
\newcommand{\ed}[1]{\mathbf{d}_{#1}}  % edge direction d_i = V[v1_i] - V[v2_i]

% -------------------------------------------------------------------------
% Steiner / triple-point
% -------------------------------------------------------------------------
\newcommand{\Sv}{\mathbf{S}}           % Steiner (Fermat-Torricelli) point
\newcommand{\Ki}{K_{i}}               % projection matrix (I-n_i n_i^T)/r_i
\newcommand{\Mmat}{M}                 % sum of K_i matrices at a triple point

% -------------------------------------------------------------------------
% Normals and unit vectors
% -------------------------------------------------------------------------
\newcommand{\nhat}{\hat{\mathbf{n}}}  % unit normal to a triangle
\newcommand{\Dhat}{\hat{\boldsymbol{\Delta}}}  % unit segment direction

% -------------------------------------------------------------------------
% Operations
% -------------------------------------------------------------------------
\newcommand{\norm}[1]{\left\|#1\right\|}
\newcommand{\abs}[1]{\left|#1\right|}
\newcommand{\ip}[2]{\langle #1,\, #2 \rangle}

% Projection perpendicular to a unit vector \uhat
\newcommand{\Proj}[1]{\bigl(\mathbf{I} - #1 #1^\top\bigr)}

% argmax operator (winner-take-all assignment in Phase 1 documents)
\DeclareMathOperator*{\argmax}{arg\,max}
\DeclareMathOperator*{\argmin}{arg\,min}

% -------------------------------------------------------------------------
% Calculus shorthands
% -------------------------------------------------------------------------
\newcommand{\pd}[2]{\dfrac{\partial #1}{\partial #2}}
\newcommand{\pdd}[3]{\dfrac{\partial^2 #1}{\partial #2\,\partial #3}}
\newcommand{\grad}{\nabla}
\newcommand{\Hess}[1]{\nabla^2 #1}

% -------------------------------------------------------------------------
% Optimization problem objects
% -------------------------------------------------------------------------
\newcommand{\Preg}{P_{\mathrm{reg}}}  % regular perimeter
\newcommand{\Pst}{P_{\mathrm{st}}}   % Steiner perimeter
\newcommand{\Areg}{A^{\mathrm{reg}}} % regular cell area
\newcommand{\Ast}{A^{\mathrm{st}}}   % Steiner area contribution
\newcommand{\Atgt}{\bar{A}}           % target area per cell
\newcommand{\Lagr}{\mathcal{L}}       % Lagrangian
\newcommand{\objfac}{\sigma}          % IPOPT objective scaling factor

% -------------------------------------------------------------------------
% Code cross-references (for "In the code..." callouts)
% -------------------------------------------------------------------------
\newcommand{\code}[1]{\texttt{#1}}
\newcommand{\codefile}[1]{\texttt{\small #1}}

% -------------------------------------------------------------------------
% Threshold dynamics / approach B (10-mbo-auction-dynamics)
% -------------------------------------------------------------------------
\newcommand{\Dlump}{D}                 % lumped mass matrix diag(v)
\newcommand{\Atau}{A_\tau}             % discrete diffusion operator (M + tau K)^{-1} M
\newcommand{\Etau}{E_\tau}             % heat-content (threshold-dynamics) energy
\newcommand{\Rcell}{R_{\mathrm{cell}}} % radius of the equal-area disc sqrt(Abar/pi)
\newcommand{\hmean}{h_{\mathrm{mean}}} % mean edge length
\newcommand{\hmax}{h_{\mathrm{max}}}   % max edge length
\newcommand{\Gt}{G_t}                  % heat kernel at time t
\newcommand{\Ktau}{K^{\mathrm{res}}_\tau} % resolvent (backward-Euler) kernel
)
% -------------------------------------------------------------------------
\usepackage{amsmath,amssymb,amsthm}
\usepackage{bm}           % \bm for bold math
\usepackage{mathtools}    % \coloneqq, dcases, etc.
\usepackage{booktabs}     % \toprule etc. in tables
\usepackage{hyperref}
\usepackage{geometry}
\usepackage{microtype}
\usepackage{parskip}
\usepackage{xcolor}
\usepackage{tcolorbox}
\usepackage{enumitem}
\usepackage{float}             % [H] float placement

\geometry{a4paper, margin=2.5cm}

% -------------------------------------------------------------------------
% Theorem-like environments
% -------------------------------------------------------------------------
\theoremstyle{definition}
\newtheorem{defn}{Definition}[section]
\newtheorem{remark}{Remark}[section]

\theoremstyle{plain}
\newtheorem{prop}{Proposition}[section]

% -------------------------------------------------------------------------
% Mesh and geometry
% -------------------------------------------------------------------------
\newcommand{\V}{\mathbf{V}}            % vertex array V[i] in R^dim
\newcommand{\F}{\mathbf{F}}            % face (triangle) array
\newcommand{\vone}[1]{v_{1,#1}}       % first vertex index of VP i
\newcommand{\vtwo}[1]{v_{2,#1}}       % second vertex index of VP i

% -------------------------------------------------------------------------
% Variable points
% -------------------------------------------------------------------------
\newcommand{\lam}{\lambda}             % scalar VP parameter
\newcommand{\Lam}{\boldsymbol{\lambda}}% full parameter vector
\newcommand{\vp}[1]{\mathbf{p}_{#1}}  % 3D position of VP i
\newcommand{\ed}[1]{\mathbf{d}_{#1}}  % edge direction d_i = V[v1_i] - V[v2_i]

% -------------------------------------------------------------------------
% Steiner / triple-point
% -------------------------------------------------------------------------
\newcommand{\Sv}{\mathbf{S}}           % Steiner (Fermat-Torricelli) point
\newcommand{\Ki}{K_{i}}               % projection matrix (I-n_i n_i^T)/r_i
\newcommand{\Mmat}{M}                 % sum of K_i matrices at a triple point

% -------------------------------------------------------------------------
% Normals and unit vectors
% -------------------------------------------------------------------------
\newcommand{\nhat}{\hat{\mathbf{n}}}  % unit normal to a triangle
\newcommand{\Dhat}{\hat{\boldsymbol{\Delta}}}  % unit segment direction

% -------------------------------------------------------------------------
% Operations
% -------------------------------------------------------------------------
\newcommand{\norm}[1]{\left\|#1\right\|}
\newcommand{\abs}[1]{\left|#1\right|}
\newcommand{\ip}[2]{\langle #1,\, #2 \rangle}

% Projection perpendicular to a unit vector \uhat
\newcommand{\Proj}[1]{\bigl(\mathbf{I} - #1 #1^\top\bigr)}

% argmax operator (winner-take-all assignment in Phase 1 documents)
\DeclareMathOperator*{\argmax}{arg\,max}
\DeclareMathOperator*{\argmin}{arg\,min}

% -------------------------------------------------------------------------
% Calculus shorthands
% -------------------------------------------------------------------------
\newcommand{\pd}[2]{\dfrac{\partial #1}{\partial #2}}
\newcommand{\pdd}[3]{\dfrac{\partial^2 #1}{\partial #2\,\partial #3}}
\newcommand{\grad}{\nabla}
\newcommand{\Hess}[1]{\nabla^2 #1}

% -------------------------------------------------------------------------
% Optimization problem objects
% -------------------------------------------------------------------------
\newcommand{\Preg}{P_{\mathrm{reg}}}  % regular perimeter
\newcommand{\Pst}{P_{\mathrm{st}}}   % Steiner perimeter
\newcommand{\Areg}{A^{\mathrm{reg}}} % regular cell area
\newcommand{\Ast}{A^{\mathrm{st}}}   % Steiner area contribution
\newcommand{\Atgt}{\bar{A}}           % target area per cell
\newcommand{\Lagr}{\mathcal{L}}       % Lagrangian
\newcommand{\objfac}{\sigma}          % IPOPT objective scaling factor

% -------------------------------------------------------------------------
% Code cross-references (for "In the code..." callouts)
% -------------------------------------------------------------------------
\newcommand{\code}[1]{\texttt{#1}}
\newcommand{\codefile}[1]{\texttt{\small #1}}

% -------------------------------------------------------------------------
% Threshold dynamics / approach B (10-mbo-auction-dynamics)
% -------------------------------------------------------------------------
\newcommand{\Dlump}{D}                 % lumped mass matrix diag(v)
\newcommand{\Atau}{A_\tau}             % discrete diffusion operator (M + tau K)^{-1} M
\newcommand{\Etau}{E_\tau}             % heat-content (threshold-dynamics) energy
\newcommand{\Rcell}{R_{\mathrm{cell}}} % radius of the equal-area disc sqrt(Abar/pi)
\newcommand{\hmean}{h_{\mathrm{mean}}} % mean edge length
\newcommand{\hmax}{h_{\mathrm{max}}}   % max edge length
\newcommand{\Gt}{G_t}                  % heat kernel at time t
\newcommand{\Ktau}{K^{\mathrm{res}}_\tau} % resolvent (backward-Euler) kernel


% -------------------------------------------------------------------------
% Packages (include here so each document only needs % shared/macros.tex
% Shared notation for all surface-partition math documents.
% Include with: % shared/macros.tex
% Shared notation for all surface-partition math documents.
% Include with: \input{../shared/macros}

% -------------------------------------------------------------------------
% Packages (include here so each document only needs \input{../shared/macros})
% -------------------------------------------------------------------------
\usepackage{amsmath,amssymb,amsthm}
\usepackage{bm}           % \bm for bold math
\usepackage{mathtools}    % \coloneqq, dcases, etc.
\usepackage{booktabs}     % \toprule etc. in tables
\usepackage{hyperref}
\usepackage{geometry}
\usepackage{microtype}
\usepackage{parskip}
\usepackage{xcolor}
\usepackage{tcolorbox}
\usepackage{enumitem}
\usepackage{float}             % [H] float placement

\geometry{a4paper, margin=2.5cm}

% -------------------------------------------------------------------------
% Theorem-like environments
% -------------------------------------------------------------------------
\theoremstyle{definition}
\newtheorem{defn}{Definition}[section]
\newtheorem{remark}{Remark}[section]

\theoremstyle{plain}
\newtheorem{prop}{Proposition}[section]

% -------------------------------------------------------------------------
% Mesh and geometry
% -------------------------------------------------------------------------
\newcommand{\V}{\mathbf{V}}            % vertex array V[i] in R^dim
\newcommand{\F}{\mathbf{F}}            % face (triangle) array
\newcommand{\vone}[1]{v_{1,#1}}       % first vertex index of VP i
\newcommand{\vtwo}[1]{v_{2,#1}}       % second vertex index of VP i

% -------------------------------------------------------------------------
% Variable points
% -------------------------------------------------------------------------
\newcommand{\lam}{\lambda}             % scalar VP parameter
\newcommand{\Lam}{\boldsymbol{\lambda}}% full parameter vector
\newcommand{\vp}[1]{\mathbf{p}_{#1}}  % 3D position of VP i
\newcommand{\ed}[1]{\mathbf{d}_{#1}}  % edge direction d_i = V[v1_i] - V[v2_i]

% -------------------------------------------------------------------------
% Steiner / triple-point
% -------------------------------------------------------------------------
\newcommand{\Sv}{\mathbf{S}}           % Steiner (Fermat-Torricelli) point
\newcommand{\Ki}{K_{i}}               % projection matrix (I-n_i n_i^T)/r_i
\newcommand{\Mmat}{M}                 % sum of K_i matrices at a triple point

% -------------------------------------------------------------------------
% Normals and unit vectors
% -------------------------------------------------------------------------
\newcommand{\nhat}{\hat{\mathbf{n}}}  % unit normal to a triangle
\newcommand{\Dhat}{\hat{\boldsymbol{\Delta}}}  % unit segment direction

% -------------------------------------------------------------------------
% Operations
% -------------------------------------------------------------------------
\newcommand{\norm}[1]{\left\|#1\right\|}
\newcommand{\abs}[1]{\left|#1\right|}
\newcommand{\ip}[2]{\langle #1,\, #2 \rangle}

% Projection perpendicular to a unit vector \uhat
\newcommand{\Proj}[1]{\bigl(\mathbf{I} - #1 #1^\top\bigr)}

% argmax operator (winner-take-all assignment in Phase 1 documents)
\DeclareMathOperator*{\argmax}{arg\,max}
\DeclareMathOperator*{\argmin}{arg\,min}

% -------------------------------------------------------------------------
% Calculus shorthands
% -------------------------------------------------------------------------
\newcommand{\pd}[2]{\dfrac{\partial #1}{\partial #2}}
\newcommand{\pdd}[3]{\dfrac{\partial^2 #1}{\partial #2\,\partial #3}}
\newcommand{\grad}{\nabla}
\newcommand{\Hess}[1]{\nabla^2 #1}

% -------------------------------------------------------------------------
% Optimization problem objects
% -------------------------------------------------------------------------
\newcommand{\Preg}{P_{\mathrm{reg}}}  % regular perimeter
\newcommand{\Pst}{P_{\mathrm{st}}}   % Steiner perimeter
\newcommand{\Areg}{A^{\mathrm{reg}}} % regular cell area
\newcommand{\Ast}{A^{\mathrm{st}}}   % Steiner area contribution
\newcommand{\Atgt}{\bar{A}}           % target area per cell
\newcommand{\Lagr}{\mathcal{L}}       % Lagrangian
\newcommand{\objfac}{\sigma}          % IPOPT objective scaling factor

% -------------------------------------------------------------------------
% Code cross-references (for "In the code..." callouts)
% -------------------------------------------------------------------------
\newcommand{\code}[1]{\texttt{#1}}
\newcommand{\codefile}[1]{\texttt{\small #1}}

% -------------------------------------------------------------------------
% Threshold dynamics / approach B (10-mbo-auction-dynamics)
% -------------------------------------------------------------------------
\newcommand{\Dlump}{D}                 % lumped mass matrix diag(v)
\newcommand{\Atau}{A_\tau}             % discrete diffusion operator (M + tau K)^{-1} M
\newcommand{\Etau}{E_\tau}             % heat-content (threshold-dynamics) energy
\newcommand{\Rcell}{R_{\mathrm{cell}}} % radius of the equal-area disc sqrt(Abar/pi)
\newcommand{\hmean}{h_{\mathrm{mean}}} % mean edge length
\newcommand{\hmax}{h_{\mathrm{max}}}   % max edge length
\newcommand{\Gt}{G_t}                  % heat kernel at time t
\newcommand{\Ktau}{K^{\mathrm{res}}_\tau} % resolvent (backward-Euler) kernel


% -------------------------------------------------------------------------
% Packages (include here so each document only needs % shared/macros.tex
% Shared notation for all surface-partition math documents.
% Include with: \input{../shared/macros}

% -------------------------------------------------------------------------
% Packages (include here so each document only needs \input{../shared/macros})
% -------------------------------------------------------------------------
\usepackage{amsmath,amssymb,amsthm}
\usepackage{bm}           % \bm for bold math
\usepackage{mathtools}    % \coloneqq, dcases, etc.
\usepackage{booktabs}     % \toprule etc. in tables
\usepackage{hyperref}
\usepackage{geometry}
\usepackage{microtype}
\usepackage{parskip}
\usepackage{xcolor}
\usepackage{tcolorbox}
\usepackage{enumitem}
\usepackage{float}             % [H] float placement

\geometry{a4paper, margin=2.5cm}

% -------------------------------------------------------------------------
% Theorem-like environments
% -------------------------------------------------------------------------
\theoremstyle{definition}
\newtheorem{defn}{Definition}[section]
\newtheorem{remark}{Remark}[section]

\theoremstyle{plain}
\newtheorem{prop}{Proposition}[section]

% -------------------------------------------------------------------------
% Mesh and geometry
% -------------------------------------------------------------------------
\newcommand{\V}{\mathbf{V}}            % vertex array V[i] in R^dim
\newcommand{\F}{\mathbf{F}}            % face (triangle) array
\newcommand{\vone}[1]{v_{1,#1}}       % first vertex index of VP i
\newcommand{\vtwo}[1]{v_{2,#1}}       % second vertex index of VP i

% -------------------------------------------------------------------------
% Variable points
% -------------------------------------------------------------------------
\newcommand{\lam}{\lambda}             % scalar VP parameter
\newcommand{\Lam}{\boldsymbol{\lambda}}% full parameter vector
\newcommand{\vp}[1]{\mathbf{p}_{#1}}  % 3D position of VP i
\newcommand{\ed}[1]{\mathbf{d}_{#1}}  % edge direction d_i = V[v1_i] - V[v2_i]

% -------------------------------------------------------------------------
% Steiner / triple-point
% -------------------------------------------------------------------------
\newcommand{\Sv}{\mathbf{S}}           % Steiner (Fermat-Torricelli) point
\newcommand{\Ki}{K_{i}}               % projection matrix (I-n_i n_i^T)/r_i
\newcommand{\Mmat}{M}                 % sum of K_i matrices at a triple point

% -------------------------------------------------------------------------
% Normals and unit vectors
% -------------------------------------------------------------------------
\newcommand{\nhat}{\hat{\mathbf{n}}}  % unit normal to a triangle
\newcommand{\Dhat}{\hat{\boldsymbol{\Delta}}}  % unit segment direction

% -------------------------------------------------------------------------
% Operations
% -------------------------------------------------------------------------
\newcommand{\norm}[1]{\left\|#1\right\|}
\newcommand{\abs}[1]{\left|#1\right|}
\newcommand{\ip}[2]{\langle #1,\, #2 \rangle}

% Projection perpendicular to a unit vector \uhat
\newcommand{\Proj}[1]{\bigl(\mathbf{I} - #1 #1^\top\bigr)}

% argmax operator (winner-take-all assignment in Phase 1 documents)
\DeclareMathOperator*{\argmax}{arg\,max}
\DeclareMathOperator*{\argmin}{arg\,min}

% -------------------------------------------------------------------------
% Calculus shorthands
% -------------------------------------------------------------------------
\newcommand{\pd}[2]{\dfrac{\partial #1}{\partial #2}}
\newcommand{\pdd}[3]{\dfrac{\partial^2 #1}{\partial #2\,\partial #3}}
\newcommand{\grad}{\nabla}
\newcommand{\Hess}[1]{\nabla^2 #1}

% -------------------------------------------------------------------------
% Optimization problem objects
% -------------------------------------------------------------------------
\newcommand{\Preg}{P_{\mathrm{reg}}}  % regular perimeter
\newcommand{\Pst}{P_{\mathrm{st}}}   % Steiner perimeter
\newcommand{\Areg}{A^{\mathrm{reg}}} % regular cell area
\newcommand{\Ast}{A^{\mathrm{st}}}   % Steiner area contribution
\newcommand{\Atgt}{\bar{A}}           % target area per cell
\newcommand{\Lagr}{\mathcal{L}}       % Lagrangian
\newcommand{\objfac}{\sigma}          % IPOPT objective scaling factor

% -------------------------------------------------------------------------
% Code cross-references (for "In the code..." callouts)
% -------------------------------------------------------------------------
\newcommand{\code}[1]{\texttt{#1}}
\newcommand{\codefile}[1]{\texttt{\small #1}}

% -------------------------------------------------------------------------
% Threshold dynamics / approach B (10-mbo-auction-dynamics)
% -------------------------------------------------------------------------
\newcommand{\Dlump}{D}                 % lumped mass matrix diag(v)
\newcommand{\Atau}{A_\tau}             % discrete diffusion operator (M + tau K)^{-1} M
\newcommand{\Etau}{E_\tau}             % heat-content (threshold-dynamics) energy
\newcommand{\Rcell}{R_{\mathrm{cell}}} % radius of the equal-area disc sqrt(Abar/pi)
\newcommand{\hmean}{h_{\mathrm{mean}}} % mean edge length
\newcommand{\hmax}{h_{\mathrm{max}}}   % max edge length
\newcommand{\Gt}{G_t}                  % heat kernel at time t
\newcommand{\Ktau}{K^{\mathrm{res}}_\tau} % resolvent (backward-Euler) kernel
)
% -------------------------------------------------------------------------
\usepackage{amsmath,amssymb,amsthm}
\usepackage{bm}           % \bm for bold math
\usepackage{mathtools}    % \coloneqq, dcases, etc.
\usepackage{booktabs}     % \toprule etc. in tables
\usepackage{hyperref}
\usepackage{geometry}
\usepackage{microtype}
\usepackage{parskip}
\usepackage{xcolor}
\usepackage{tcolorbox}
\usepackage{enumitem}
\usepackage{float}             % [H] float placement

\geometry{a4paper, margin=2.5cm}

% -------------------------------------------------------------------------
% Theorem-like environments
% -------------------------------------------------------------------------
\theoremstyle{definition}
\newtheorem{defn}{Definition}[section]
\newtheorem{remark}{Remark}[section]

\theoremstyle{plain}
\newtheorem{prop}{Proposition}[section]

% -------------------------------------------------------------------------
% Mesh and geometry
% -------------------------------------------------------------------------
\newcommand{\V}{\mathbf{V}}            % vertex array V[i] in R^dim
\newcommand{\F}{\mathbf{F}}            % face (triangle) array
\newcommand{\vone}[1]{v_{1,#1}}       % first vertex index of VP i
\newcommand{\vtwo}[1]{v_{2,#1}}       % second vertex index of VP i

% -------------------------------------------------------------------------
% Variable points
% -------------------------------------------------------------------------
\newcommand{\lam}{\lambda}             % scalar VP parameter
\newcommand{\Lam}{\boldsymbol{\lambda}}% full parameter vector
\newcommand{\vp}[1]{\mathbf{p}_{#1}}  % 3D position of VP i
\newcommand{\ed}[1]{\mathbf{d}_{#1}}  % edge direction d_i = V[v1_i] - V[v2_i]

% -------------------------------------------------------------------------
% Steiner / triple-point
% -------------------------------------------------------------------------
\newcommand{\Sv}{\mathbf{S}}           % Steiner (Fermat-Torricelli) point
\newcommand{\Ki}{K_{i}}               % projection matrix (I-n_i n_i^T)/r_i
\newcommand{\Mmat}{M}                 % sum of K_i matrices at a triple point

% -------------------------------------------------------------------------
% Normals and unit vectors
% -------------------------------------------------------------------------
\newcommand{\nhat}{\hat{\mathbf{n}}}  % unit normal to a triangle
\newcommand{\Dhat}{\hat{\boldsymbol{\Delta}}}  % unit segment direction

% -------------------------------------------------------------------------
% Operations
% -------------------------------------------------------------------------
\newcommand{\norm}[1]{\left\|#1\right\|}
\newcommand{\abs}[1]{\left|#1\right|}
\newcommand{\ip}[2]{\langle #1,\, #2 \rangle}

% Projection perpendicular to a unit vector \uhat
\newcommand{\Proj}[1]{\bigl(\mathbf{I} - #1 #1^\top\bigr)}

% argmax operator (winner-take-all assignment in Phase 1 documents)
\DeclareMathOperator*{\argmax}{arg\,max}
\DeclareMathOperator*{\argmin}{arg\,min}

% -------------------------------------------------------------------------
% Calculus shorthands
% -------------------------------------------------------------------------
\newcommand{\pd}[2]{\dfrac{\partial #1}{\partial #2}}
\newcommand{\pdd}[3]{\dfrac{\partial^2 #1}{\partial #2\,\partial #3}}
\newcommand{\grad}{\nabla}
\newcommand{\Hess}[1]{\nabla^2 #1}

% -------------------------------------------------------------------------
% Optimization problem objects
% -------------------------------------------------------------------------
\newcommand{\Preg}{P_{\mathrm{reg}}}  % regular perimeter
\newcommand{\Pst}{P_{\mathrm{st}}}   % Steiner perimeter
\newcommand{\Areg}{A^{\mathrm{reg}}} % regular cell area
\newcommand{\Ast}{A^{\mathrm{st}}}   % Steiner area contribution
\newcommand{\Atgt}{\bar{A}}           % target area per cell
\newcommand{\Lagr}{\mathcal{L}}       % Lagrangian
\newcommand{\objfac}{\sigma}          % IPOPT objective scaling factor

% -------------------------------------------------------------------------
% Code cross-references (for "In the code..." callouts)
% -------------------------------------------------------------------------
\newcommand{\code}[1]{\texttt{#1}}
\newcommand{\codefile}[1]{\texttt{\small #1}}

% -------------------------------------------------------------------------
% Threshold dynamics / approach B (10-mbo-auction-dynamics)
% -------------------------------------------------------------------------
\newcommand{\Dlump}{D}                 % lumped mass matrix diag(v)
\newcommand{\Atau}{A_\tau}             % discrete diffusion operator (M + tau K)^{-1} M
\newcommand{\Etau}{E_\tau}             % heat-content (threshold-dynamics) energy
\newcommand{\Rcell}{R_{\mathrm{cell}}} % radius of the equal-area disc sqrt(Abar/pi)
\newcommand{\hmean}{h_{\mathrm{mean}}} % mean edge length
\newcommand{\hmax}{h_{\mathrm{max}}}   % max edge length
\newcommand{\Gt}{G_t}                  % heat kernel at time t
\newcommand{\Ktau}{K^{\mathrm{res}}_\tau} % resolvent (backward-Euler) kernel
)
% -------------------------------------------------------------------------
\usepackage{amsmath,amssymb,amsthm}
\usepackage{bm}           % \bm for bold math
\usepackage{mathtools}    % \coloneqq, dcases, etc.
\usepackage{booktabs}     % \toprule etc. in tables
\usepackage{hyperref}
\usepackage{geometry}
\usepackage{microtype}
\usepackage{parskip}
\usepackage{xcolor}
\usepackage{tcolorbox}
\usepackage{enumitem}
\usepackage{float}             % [H] float placement

\geometry{a4paper, margin=2.5cm}

% -------------------------------------------------------------------------
% Theorem-like environments
% -------------------------------------------------------------------------
\theoremstyle{definition}
\newtheorem{defn}{Definition}[section]
\newtheorem{remark}{Remark}[section]

\theoremstyle{plain}
\newtheorem{prop}{Proposition}[section]

% -------------------------------------------------------------------------
% Mesh and geometry
% -------------------------------------------------------------------------
\newcommand{\V}{\mathbf{V}}            % vertex array V[i] in R^dim
\newcommand{\F}{\mathbf{F}}            % face (triangle) array
\newcommand{\vone}[1]{v_{1,#1}}       % first vertex index of VP i
\newcommand{\vtwo}[1]{v_{2,#1}}       % second vertex index of VP i

% -------------------------------------------------------------------------
% Variable points
% -------------------------------------------------------------------------
\newcommand{\lam}{\lambda}             % scalar VP parameter
\newcommand{\Lam}{\boldsymbol{\lambda}}% full parameter vector
\newcommand{\vp}[1]{\mathbf{p}_{#1}}  % 3D position of VP i
\newcommand{\ed}[1]{\mathbf{d}_{#1}}  % edge direction d_i = V[v1_i] - V[v2_i]

% -------------------------------------------------------------------------
% Steiner / triple-point
% -------------------------------------------------------------------------
\newcommand{\Sv}{\mathbf{S}}           % Steiner (Fermat-Torricelli) point
\newcommand{\Ki}{K_{i}}               % projection matrix (I-n_i n_i^T)/r_i
\newcommand{\Mmat}{M}                 % sum of K_i matrices at a triple point

% -------------------------------------------------------------------------
% Normals and unit vectors
% -------------------------------------------------------------------------
\newcommand{\nhat}{\hat{\mathbf{n}}}  % unit normal to a triangle
\newcommand{\Dhat}{\hat{\boldsymbol{\Delta}}}  % unit segment direction

% -------------------------------------------------------------------------
% Operations
% -------------------------------------------------------------------------
\newcommand{\norm}[1]{\left\|#1\right\|}
\newcommand{\abs}[1]{\left|#1\right|}
\newcommand{\ip}[2]{\langle #1,\, #2 \rangle}

% Projection perpendicular to a unit vector \uhat
\newcommand{\Proj}[1]{\bigl(\mathbf{I} - #1 #1^\top\bigr)}

% argmax operator (winner-take-all assignment in Phase 1 documents)
\DeclareMathOperator*{\argmax}{arg\,max}
\DeclareMathOperator*{\argmin}{arg\,min}

% -------------------------------------------------------------------------
% Calculus shorthands
% -------------------------------------------------------------------------
\newcommand{\pd}[2]{\dfrac{\partial #1}{\partial #2}}
\newcommand{\pdd}[3]{\dfrac{\partial^2 #1}{\partial #2\,\partial #3}}
\newcommand{\grad}{\nabla}
\newcommand{\Hess}[1]{\nabla^2 #1}

% -------------------------------------------------------------------------
% Optimization problem objects
% -------------------------------------------------------------------------
\newcommand{\Preg}{P_{\mathrm{reg}}}  % regular perimeter
\newcommand{\Pst}{P_{\mathrm{st}}}   % Steiner perimeter
\newcommand{\Areg}{A^{\mathrm{reg}}} % regular cell area
\newcommand{\Ast}{A^{\mathrm{st}}}   % Steiner area contribution
\newcommand{\Atgt}{\bar{A}}           % target area per cell
\newcommand{\Lagr}{\mathcal{L}}       % Lagrangian
\newcommand{\objfac}{\sigma}          % IPOPT objective scaling factor

% -------------------------------------------------------------------------
% Code cross-references (for "In the code..." callouts)
% -------------------------------------------------------------------------
\newcommand{\code}[1]{\texttt{#1}}
\newcommand{\codefile}[1]{\texttt{\small #1}}

% -------------------------------------------------------------------------
% Threshold dynamics / approach B (10-mbo-auction-dynamics)
% -------------------------------------------------------------------------
\newcommand{\Dlump}{D}                 % lumped mass matrix diag(v)
\newcommand{\Atau}{A_\tau}             % discrete diffusion operator (M + tau K)^{-1} M
\newcommand{\Etau}{E_\tau}             % heat-content (threshold-dynamics) energy
\newcommand{\Rcell}{R_{\mathrm{cell}}} % radius of the equal-area disc sqrt(Abar/pi)
\newcommand{\hmean}{h_{\mathrm{mean}}} % mean edge length
\newcommand{\hmax}{h_{\mathrm{max}}}   % max edge length
\newcommand{\Gt}{G_t}                  % heat kernel at time t
\newcommand{\Ktau}{K^{\mathrm{res}}_\tau} % resolvent (backward-Euler) kernel


\hypersetup{
  colorlinks = true,
  linkcolor  = blue!60!black,
  citecolor  = green!50!black,
  urlcolor   = blue!60!black,
  pdftitle   = {Phase 2 Optimization Derivatives},
  pdfauthor  = {surface-partition project}
}

% =========================================================================
\title{%
  \textbf{Phase~2 Optimization Derivatives}\\[0.5em]
  \large Perimeter and Area Quantities Implemented in \codefile{surface-partition}
}
\author{}
\date{May 2026}
% =========================================================================

\begin{document}

\maketitle

\begin{abstract}
This document derives every quantity that the Phase~2 perimeter-refinement
optimizer evaluates during an IPOPT iteration: the perimeter objective
$P(\Lam)$, its gradient $\grad P$, and its Hessian $\Hess{P}$; the area
constraint values $A_k(\Lam)$, their Jacobian $J$, and the
multiplier-weighted Hessian $\sum_k \mu_k \Hess{A_k}$.  Quantities
belonging to \emph{regular} boundary segments and triangles are derived in
closed form.  Quantities involving \emph{Steiner / triple-point} geometry
have analytical forward values; their analytical first and second
derivatives are derived in the companion document
\codefile{docs/math/03-analytical-steiner-derivatives}.  All formulas are
cross-referenced to the implementing Python functions in
\codefile{src/partition/}.
\end{abstract}

\tableofcontents
\newpage

% =========================================================================
\section{Overview}
% =========================================================================

Phase~2 of the \codefile{surface-partition} pipeline solves the constrained
nonlinear program
\begin{equation}
  \min_{\Lam \in \mathbb{R}^n}\; P(\Lam)
  \quad\text{subject to}\quad
  A_k(\Lam) = \Atgt,\quad k = 0,\ldots,N-2,
  \label{eq:nlp}
\end{equation}
where
\begin{itemize}[noitemsep]
  \item $\Lam = (\lam_1,\ldots,\lam_n)^\top$ collects the scalar parameters
        of all active Variable Points (VPs),
  \item $P(\Lam)$ is the total boundary perimeter,
  \item $A_k(\Lam)$ is the surface area of cell $k$, and
  \item $\Atgt = A_{\mathrm{total}}/N$ is the equal-area target.
\end{itemize}
The solver used is IPOPT~\cite{wachter2006ipopt} (Interior-Point OPTimizer).
At each iteration IPOPT requires the Lagrangian Hessian
\begin{equation}
  H = \objfac\,\Hess{P} + \sum_{k=0}^{N-2} \mu_k\,\Hess{A_k},
  \label{eq:lag-hess}
\end{equation}
in addition to gradients and Jacobians.  The scalar $\objfac$ (objective
scaling factor) and the Lagrange multipliers $\mu_k$ are supplied by IPOPT
at each callback.

\paragraph{Structure of this document.}
Sections~\ref{sec:notation}--\ref{sec:area} cover the regular (non-Steiner)
path, which is fully analytical.  Section~\ref{sec:steiner} covers the
Steiner / triple-point path; its forward values are analytical and its
derivatives, also analytical, are derived in the companion document
\codefile{docs/math/03-analytical-steiner-derivatives}.
Section~\ref{sec:assembly} describes how the pieces assemble into the
Lagrangian Hessian and its sparsity pattern.

% =========================================================================
\section{Notation and Setup}
\label{sec:notation}
% =========================================================================

\subsection{The mesh}

The surface is represented as a triangulated mesh $(\V, \F)$, where
$\V \in \mathbb{R}^{M \times 3}$ is the vertex position matrix and $\F$
is the face array.  The mesh lives in $\mathbb{R}^3$; all cross-product
formulas are three-dimensional unless otherwise noted.

\subsection{Variable Points and the $\lambda$ convention}
\label{sec:lambda-convention}

A \emph{Variable Point} (VP) $i$ sits on a mesh edge $(v_{1,i},\, v_{2,i})$
with the convention $v_{1,i} < v_{2,i}$ (smaller index first).  Its
three-dimensional position is the affine combination
\begin{equation}
  \vp{i}(\lam_i)
    \coloneqq \lam_i\, \V[v_{1,i}] + (1 - \lam_i)\, \V[v_{2,i}],
  \qquad \lam_i \in [0,1].
  \label{eq:vp-position}
\end{equation}
\begin{itemize}[noitemsep]
  \item $\lam_i = 1$: VP is at the vertex with the \emph{smaller} index.
  \item $\lam_i = 0$: VP is at the vertex with the \emph{larger} index.
\end{itemize}

\begin{tcolorbox}[colback=blue!4!white, colframe=blue!40!black,
                  title={Code: VP position}, fontupper=\small]
\codefile{vectorized\_perimeter.py}:
\code{\_compute\_vp\_positions(pa)}\\[2pt]
\texttt{p[i] = pa.vp\_lambda[i] * V[pa.vp\_edge\_v1[i]]\\
\phantom{p[i] = }+ (1 - pa.vp\_lambda[i]) * V[pa.vp\_edge\_v2[i]]}
\end{tcolorbox}

\subsection{The edge direction identity}
\label{sec:edge-direction}

Define the \emph{edge direction} of VP $i$:
\begin{equation}
  \ed{i} \coloneqq \V[v_{1,i}] - \V[v_{2,i}] \in \mathbb{R}^3.
  \label{eq:edge-dir}
\end{equation}
Differentiating \eqref{eq:vp-position} with respect to $\lam_i$ gives the
fundamental identity used throughout all derivative calculations:
\begin{equation}
  \boxed{\pd{\vp{i}}{\lam_i} = \ed{i}, \qquad
  \pdd{\vp{i}}{\lam_i}{\lam_i} = \mathbf{0}.}
  \label{eq:vp-deriv}
\end{equation}
The position is \emph{affine} in $\lam_i$, so all second derivatives of
$\vp{i}$ vanish.  This dramatically simplifies the Hessian calculations
in Sections~\ref{sec:perimeter} and~\ref{sec:area}: every second-order
term that would involve $\partial^2 \vp{i}/\partial\lam_i^2$ is zero.

% =========================================================================
\section{Regular Perimeter}
\label{sec:perimeter}
% =========================================================================

\subsection{Perimeter value}
\label{sec:perimeter-value}

The boundary between cells is represented as a collection of
\emph{segments}, each connecting two VPs.  Let segment $s$ connect VP $a$
and VP $b$, and let $w_s$ be its multiplicity weight:
\begin{equation}
  w_s = \begin{cases} 2 & \text{if the segment separates two distinct cells
                            ($\mathrm{cell}_a \neq \mathrm{cell}_b$),} \\
                       1 & \text{otherwise (fallback: triple-point void
                            segment, $\mathrm{cell}_a = \mathrm{cell}_b$).}
        \end{cases}
  \label{eq:weight}
\end{equation}
The weight $w_s = 2$ reflects the double-counting convention: each proper
inter-cell boundary edge is shared by exactly two cells and therefore
counted once in each cell's perimeter.  The total regular perimeter is
\begin{equation}
  \Preg(\Lam)
    = \sum_s w_s\, \norm{\vp{a(s)} - \vp{b(s)}}.
  \label{eq:perimeter}
\end{equation}

\begin{tcolorbox}[colback=blue!4!white, colframe=blue!40!black,
                  title={Code: regular perimeter value}, fontupper=\small]
\codefile{vectorized\_perimeter.py}: \code{compute\_total\_perimeter(pa)}
\end{tcolorbox}

\subsection{Perimeter gradient}
\label{sec:perimeter-gradient}

For a single segment $s$ with endpoints $\vp{a}$ and $\vp{b}$, define
\begin{equation}
  \boldsymbol{\Delta}_s \coloneqq \vp{a} - \vp{b}, \qquad
  r_s \coloneqq \norm{\boldsymbol{\Delta}_s}, \qquad
  \Dhat_s \coloneqq \boldsymbol{\Delta}_s / r_s.
\end{equation}
The contribution of $s$ to the perimeter is $L_s = w_s r_s$.  Applying the
chain rule with identity \eqref{eq:vp-deriv}:
\begin{align}
  \pd{L_s}{\lam_a}
    &= w_s \, \frac{\boldsymbol{\Delta}_s \cdot \ed{a}}{r_s},
  \label{eq:dL-da} \\
  \pd{L_s}{\lam_b}
    &= w_s \, \frac{-\boldsymbol{\Delta}_s \cdot \ed{b}}{r_s}.
  \label{eq:dL-db}
\end{align}
The sign difference arises because $\partial\boldsymbol{\Delta}_s/\partial\lam_a =
\ed{a}$ but $\partial\boldsymbol{\Delta}_s/\partial\lam_b = -\ed{b}$.

The total gradient accumulates over all segments:
\begin{equation}
  \pd{\Preg}{\lam_i}
    = \sum_{\substack{s:\, a(s)=i}} w_s\,
        \frac{\boldsymbol{\Delta}_s \cdot \ed{i}}{r_s}
    + \sum_{\substack{s:\, b(s)=i}} w_s\,
        \frac{-\boldsymbol{\Delta}_s \cdot \ed{i}}{r_s}.
  \label{eq:grad-P}
\end{equation}

\begin{tcolorbox}[colback=blue!4!white, colframe=blue!40!black,
                  title={Code: perimeter gradient}, fontupper=\small]
\codefile{vectorized\_perimeter.py}: \code{compute\_perimeter\_gradient(pa)}\\[2pt]
\texttt{grad\_vp1[s] = sum(diff * dv1, axis=1) / lengths[s]}\\
\texttt{grad\_vp2[s] = sum(-diff * dv2, axis=1) / lengths[s]}
\end{tcolorbox}

\subsection{Perimeter Hessian}
\label{sec:perimeter-hessian}

We compute $\partial^2 L_s / \partial\lam_i\,\partial\lam_j$ for the same
single segment.  The key tool is the Hessian of the Euclidean norm
$g(\boldsymbol{\Delta}) = \norm{\boldsymbol{\Delta}}$:
\begin{equation}
  \frac{\partial^2 \norm{\boldsymbol{\Delta}}}{\partial \Delta_m\,\partial \Delta_n}
  = \frac{\delta_{mn}}{r} - \frac{\Delta_m \Delta_n}{r^3}
  = \frac{1}{r}\Proj{\Dhat}_{mn},
  \label{eq:norm-hessian}
\end{equation}
where $\Proj{\Dhat} = \mathbf{I} - \Dhat\Dhat^\top$ is the projection onto
the subspace perpendicular to $\boldsymbol{\Delta}$.

Since $\partial\boldsymbol{\Delta}/\partial\lam_a = \ed{a}$ and
$\partial\boldsymbol{\Delta}/\partial\lam_b = -\ed{b}$, the chain rule gives:
\begin{align}
  H^{(s)}_{aa}
    &= w_s\,\ed{a}^\top \,\frac{\Proj{\Dhat_s}}{r_s}\, \ed{a}
     = \frac{w_s}{r_s}\!\left(\norm{\ed{a}}^2 - (\ed{a}\cdot\Dhat_s)^2\right),
  \label{eq:H_aa} \\[4pt]
  H^{(s)}_{bb}
    &= w_s\,\ed{b}^\top \,\frac{\Proj{\Dhat_s}}{r_s}\, \ed{b}
     = \frac{w_s}{r_s}\!\left(\norm{\ed{b}}^2 - (\ed{b}\cdot\Dhat_s)^2\right),
  \label{eq:H_bb} \\[4pt]
  H^{(s)}_{ab}
    &= w_s\,\ed{a}^\top \,\frac{\Proj{\Dhat_s}}{r_s}\, (-\ed{b})
     = -\frac{w_s}{r_s}\!\left(\ed{a}\cdot\ed{b}
         - (\ed{a}\cdot\Dhat_s)(\ed{b}\cdot\Dhat_s)\right).
  \label{eq:H_ab}
\end{align}

\begin{remark}
The factor $\Proj{\Dhat_s}/r_s$ is the projection perpendicular to the
segment direction, divided by the segment length.  When $\ed{i}$ is
\emph{parallel} to $\boldsymbol{\Delta}_s$, the entry is zero — moving
along the segment does not change its length to first order, which the
gradient already captures.  When $\ed{i}$ is \emph{perpendicular}, the
entry equals $w_s\norm{\ed{i}}^2/r_s$ (maximum curvature).
\end{remark}

The full perimeter Hessian is the sum over all segments:
\begin{equation}
  \pdd{\Preg}{\lam_i}{\lam_j}
    = \sum_s H^{(s)}_{ij}\,\mathbf{1}[\{i,j\}\subseteq\{a(s),b(s)\}].
  \label{eq:hess-P}
\end{equation}
Only pairs $(i,j)$ that are both endpoints of a common segment contribute;
the pattern is determined at compile time from \code{pa.seg\_vp1} and
\code{pa.seg\_vp2}.

\begin{tcolorbox}[colback=blue!4!white, colframe=blue!40!black,
                  title={Code: perimeter Hessian}, fontupper=\small]
\codefile{vectorized\_perimeter.py}: \code{compute\_perimeter\_hessian\_sparse(pa)}\\[2pt]
\texttt{H\_aa = w * (da\_dot\_da - da\_dot\_dh**2) / r}\\
\texttt{H\_bb = w * (db\_dot\_db - db\_dot\_dh**2) / r}\\
\texttt{H\_ab = -w * (da\_dot\_db - da\_dot\_dh * db\_dot\_dh) / r}\\
Accumulated with \texttt{np.add.at} into the pre-allocated sparse values
array at offsets \code{pa.seg\_hess\_off\_aa}, \code{pa.seg\_hess\_off\_bb},
\code{pa.seg\_hess\_off\_ab}.
\end{tcolorbox}

% =========================================================================
\section{Regular Cell Area}
\label{sec:area}
% =========================================================================

A mesh triangle $T$ is a \emph{boundary triangle} if the cell boundary
crosses it — that is, if its three vertices do not all belong to the same
cell.  Each boundary triangle contributes to the area of exactly one cell.
There are two structural cases depending on how many of the triangle's
mesh vertices lie \emph{inside} the cell (Fig.~\ref{fig:btri}).

\begin{figure}[h]
\centering
\begin{tabular}{cc}
\begin{minipage}{0.45\textwidth}
\centering
\textit{1-inside:}\\[4pt]
Triangle $T = (v_\mathrm{in},\, v_\mathrm{out1},\, v_\mathrm{out2})$.\\
VPs: $\mathrm{pc}_1$ on edge $(v_\mathrm{in}, v_\mathrm{out1})$,\\
\phantom{VPs: }$\mathrm{pc}_2$ on edge $(v_\mathrm{in}, v_\mathrm{out2})$.\\[4pt]
Cell fragment $= \triangle(p_\mathrm{in},\, \mathrm{pc}_1,\, \mathrm{pc}_2)$.
\end{minipage}
&
\begin{minipage}{0.45\textwidth}
\centering
\textit{2-inside:}\\[4pt]
Triangle $T = (v_{\mathrm{in}1},\, v_{\mathrm{in}2},\, v_\mathrm{out})$.\\
VPs: $\mathrm{pc}_1$ on edge $(v_{\mathrm{in}1}, v_\mathrm{out})$,\\
\phantom{VPs: }$\mathrm{pc}_2$ on edge $(v_{\mathrm{in}2}, v_\mathrm{out})$.\\[4pt]
Cell fragment $= \mathrm{quad}(p_{\mathrm{in}1},\, \mathrm{pc}_1,\,
                               \mathrm{pc}_2,\, p_{\mathrm{in}2})$.
\end{minipage}
\end{tabular}
\caption{The two boundary-triangle cases.  Filled circles mark vertices
inside the cell; open circles mark vertices outside.  Stars mark Variable
Points (VPs) on the crossing edges.}
\label{fig:btri}
\end{figure}

\begin{tcolorbox}[colback=blue!4!white, colframe=blue!40!black,
                  title={Code: area arrays}, fontupper=\small]
Arrays \code{pa.btri\_n\_inside}, \code{pa.btri\_v\_in}, \code{pa.btri\_vp1},
\code{pa.btri\_vp2}, \code{pa.btri\_cell} fully characterize each boundary
triangle.
\end{tcolorbox}

\subsection{1-inside case}
\label{sec:area-1inside}

Let $p_\mathrm{in} = \V[v_\mathrm{in}]$ (fixed mesh vertex),
$\mathrm{pc}_1 = \vp{1}(\lam_1)$, $\mathrm{pc}_2 = \vp{2}(\lam_2)$.
Define
\begin{equation}
  \mathbf{u} \coloneqq \mathrm{pc}_1 - p_\mathrm{in},\qquad
  \mathbf{v} \coloneqq \mathrm{pc}_2 - p_\mathrm{in}.
\end{equation}

\subsubsection{Area value}

\begin{equation}
  A = \tfrac{1}{2}\norm{\mathbf{u} \times \mathbf{v}}.
  \label{eq:area-1inside}
\end{equation}
Let $\mathbf{C} \coloneqq \mathbf{u} \times \mathbf{v}$,
$\norm{\mathbf{C}} = 2A$, and $\nhat \coloneqq \mathbf{C}/\norm{\mathbf{C}}$.

\subsubsection{Area Jacobian}

Differentiating \eqref{eq:area-1inside} with respect to $\lam_1$ and
$\lam_2$, and using $\partial\mathbf{u}/\partial\lam_1 = \ed{1}$,
$\partial\mathbf{v}/\partial\lam_1 = \mathbf{0}$,
$\partial\mathbf{u}/\partial\lam_2 = \mathbf{0}$,
$\partial\mathbf{v}/\partial\lam_2 = \ed{2}$:
\begin{equation}
  \pd{A}{\lam_1}
    = \tfrac{1}{2}\,\nhat \cdot (\ed{1} \times \mathbf{v}),
  \qquad
  \pd{A}{\lam_2}
    = \tfrac{1}{2}\,\nhat \cdot (\mathbf{u} \times \ed{2}).
  \label{eq:area-jac-1inside}
\end{equation}

\begin{tcolorbox}[colback=blue!4!white, colframe=blue!40!black,
                  title={Code: 1-inside Jacobian}, fontupper=\small]
\codefile{vectorized\_area.py}: \code{compute\_area\_jacobian\_analytical},
1-inside branch:\\[2pt]
\texttt{dA\_dl1 = 0.5 * np.sum(n\_hat * np.cross(d1, v), axis=1)}\\
\texttt{dA\_dl2 = 0.5 * np.sum(n\_hat * np.cross(u, d2), axis=1)}
\end{tcolorbox}

\subsubsection{Area Hessian}
\label{sec:area-hessian-1inside}

Introduce the auxiliary cross products
\begin{equation}
  \mathbf{g}_1 \coloneqq \ed{1} \times \mathbf{v}
    = \pd{\mathbf{C}}{\lam_1},
  \qquad
  \mathbf{g}_2 \coloneqq \mathbf{u} \times \ed{2}
    = \pd{\mathbf{C}}{\lam_2}.
\end{equation}
Differentiating the unit-normal identity $A = \tfrac{1}{2}\norm{\mathbf{C}}$
a second time, using the standard formula for the derivative of a unit
vector, and noting that $\partial^2 \mathbf{C}/\partial\lam_1^2 =
\partial^2\mathbf{C}/\partial\lam_2^2 = \mathbf{0}$ (because
$\partial^2\vp{i}/\partial\lam_i^2 = \mathbf{0}$):
\begin{align}
  \pdd{A}{\lam_1}{\lam_1}
    &= \frac{\norm{\mathbf{g}_1}^2 - (\nhat\cdot\mathbf{g}_1)^2}
            {2\norm{\mathbf{C}}},
  \label{eq:H11-1inside} \\[4pt]
  \pdd{A}{\lam_2}{\lam_2}
    &= \frac{\norm{\mathbf{g}_2}^2 - (\nhat\cdot\mathbf{g}_2)^2}
            {2\norm{\mathbf{C}}},
  \label{eq:H22-1inside} \\[4pt]
  \pdd{A}{\lam_1}{\lam_2}
    &= \frac{\mathbf{g}_1\cdot\mathbf{g}_2
             - (\nhat\cdot\mathbf{g}_1)(\nhat\cdot\mathbf{g}_2)}
            {2\norm{\mathbf{C}}}
       + \frac{\nhat \cdot (\ed{1} \times \ed{2})}{2}.
  \label{eq:H12-1inside}
\end{align}

The off-diagonal term \eqref{eq:H12-1inside} has two parts:
\begin{itemize}[noitemsep]
  \item The first fraction arises from differentiating $\nhat$ — it
        projects $\mathbf{g}_1$ and $\mathbf{g}_2$ perpendicular to $\nhat$.
  \item The second term $\nhat\cdot(\ed{1}\times\ed{2})/2$ arises from
        $\partial\mathbf{g}_1/\partial\lam_2 = \ed{1}\times\ed{2}$
        (differentiating $\mathbf{g}_1 = \ed{1}\times\mathbf{v}$ with
        respect to $\lam_2$, since $\partial\mathbf{v}/\partial\lam_2 = \ed{2}$).
\end{itemize}

\begin{tcolorbox}[colback=blue!4!white, colframe=blue!40!black,
                  title={Code: 1-inside Hessian}, fontupper=\small]
\codefile{vectorized\_area.py}: \code{compute\_area\_hessian\_sparse},
1-inside branch:\\[2pt]
\texttt{g1 = np.cross(d1, v) \# $\partial C/\partial\lambda_1$}\\
\texttt{g2 = np.cross(u, d2) \# $\partial C/\partial\lambda_2$}\\
\texttt{d1xd2 = np.cross(d1, d2)}\\
\texttt{H\_11 = (g1\_dot\_g1 - n\_dot\_g1**2) / (2.0 * normC)}\\
\texttt{H\_22 = (g2\_dot\_g2 - n\_dot\_g2**2) / (2.0 * normC)}\\
\texttt{H\_12 = ((g1\_dot\_g2 - n\_dot\_g1 * n\_dot\_g2) / normC}\\
\texttt{\phantom{H\_12 = (}+ n\_dot\_d1xd2) / 2.0}
\end{tcolorbox}

\subsection{2-inside case}
\label{sec:area-2inside}

Let $p_{\mathrm{in}1} = \V[v_{\mathrm{in}1}]$,
$p_{\mathrm{in}2} = \V[v_{\mathrm{in}2}]$ (fixed mesh vertices),
$\mathrm{pc}_1 = \vp{1}(\lam_1)$, $\mathrm{pc}_2 = \vp{2}(\lam_2)$.

\subsubsection{Area value}

The cell's contribution is the quadrilateral area, split into two
sub-triangles:
\begin{equation}
  A = A_1 + A_2,
\end{equation}
where
\begin{align}
  A_1 &= \tfrac{1}{2}\norm{(\mathrm{pc}_1 - p_{\mathrm{in}1})
                   \times (\mathrm{pc}_2 - p_{\mathrm{in}1})},
  \label{eq:A1-2inside} \\
  A_2 &= \tfrac{1}{2}\norm{(\mathrm{pc}_2 - p_{\mathrm{in}1})
                   \times (p_{\mathrm{in}2} - p_{\mathrm{in}1})}.
  \label{eq:A2-2inside}
\end{align}

Introduce the sub-triangle vectors:
\begin{alignat}{2}
  \mathbf{u}_1 &= \mathrm{pc}_1 - p_{\mathrm{in}1}, &\quad
  \mathbf{v}_1 &= \mathrm{pc}_2 - p_{\mathrm{in}1}, \\
  \mathbf{u}_2 &= \mathrm{pc}_2 - p_{\mathrm{in}1}, &\quad
  \mathbf{v}_2 &= p_{\mathrm{in}2} - p_{\mathrm{in}1}.
\end{alignat}
Note that $\mathbf{v}_2$ is \emph{constant} (both vertices are fixed).

\subsubsection{Area Jacobian}

VP~1 appears only in $A_1$ (through $\mathrm{pc}_1 = \mathbf{u}_1 +
p_{\mathrm{in}1}$).  VP~2 appears in both $A_1$ and $A_2$ (through
$\mathrm{pc}_2$):
\begin{align}
  \pd{A}{\lam_1}
    &= \pd{A_1}{\lam_1}
     = \tfrac{1}{2}\,\hat{\mathbf{n}}_1 \cdot (\ed{1} \times \mathbf{v}_1),
  \label{eq:jac2-1} \\[4pt]
  \pd{A}{\lam_2}
    &= \pd{A_1}{\lam_2} + \pd{A_2}{\lam_2}
     = \tfrac{1}{2}\,\hat{\mathbf{n}}_1 \cdot (\mathbf{u}_1 \times \ed{2})
       + \tfrac{1}{2}\,\hat{\mathbf{n}}_2 \cdot (\ed{2} \times \mathbf{v}_2),
  \label{eq:jac2-2}
\end{align}
where $\hat{\mathbf{n}}_1$, $\hat{\mathbf{n}}_2$ are the unit normals of
sub-triangles 1 and 2 respectively.

\subsubsection{Area Hessian}

Define auxiliary cross products for each sub-triangle following the same
pattern as Section~\ref{sec:area-hessian-1inside}:
\begin{alignat}{2}
  \mathbf{g}_{1,1} &= \ed{1} \times \mathbf{v}_1, &\qquad
  \mathbf{g}_{2,1} &= \mathbf{u}_1 \times \ed{2},
  \quad\text{(sub-triangle 1)} \\
  \mathbf{g}_{2,2} &= \ed{2} \times \mathbf{v}_2.
  \qquad\text{(sub-triangle 2)}
\end{alignat}

Sub-triangle 2 does not involve $\lam_1$, so its only contribution to the
Hessian is to $\partial^2 A / \partial\lam_2^2$.  The full Hessian entries
are:
\begin{align}
  \pdd{A}{\lam_1}{\lam_1} &= H^{(1)}_{11},
  \label{eq:H11-2inside} \\[4pt]
  \pdd{A}{\lam_2}{\lam_2} &= H^{(1)}_{22} + H^{(2)}_{22},
  \label{eq:H22-2inside} \\[4pt]
  \pdd{A}{\lam_1}{\lam_2} &= H^{(1)}_{12},
  \label{eq:H12-2inside}
\end{align}
where $H^{(j)}_{mn}$ denotes the Hessian entry from sub-triangle $j$,
computed by the formulas \eqref{eq:H11-1inside}--\eqref{eq:H12-1inside}
with the appropriate sub-triangle vectors substituted.  Explicitly:
\begin{align}
  H^{(2)}_{22}
    &= \frac{\norm{\mathbf{g}_{2,2}}^2
             - (\hat{\mathbf{n}}_2\cdot\mathbf{g}_{2,2})^2}
            {2\norm{\mathbf{C}_2}},
  \label{eq:H22-sub2}
\end{align}
where $\mathbf{C}_2 = \mathbf{u}_2 \times \mathbf{v}_2$ and $\norm{\mathbf{C}_2}
= 2A_2$.

\begin{tcolorbox}[colback=blue!4!white, colframe=blue!40!black,
                  title={Code: 2-inside Hessian}, fontupper=\small]
\codefile{vectorized\_area.py}: \code{compute\_area\_hessian\_sparse},
2-inside branch:\\[2pt]
\texttt{H\_11 = H\_11\_sub1}\\
\texttt{H\_22 = H\_22\_sub1 + H\_22\_sub2}\\
\texttt{H\_12 = H\_12\_sub1}
\end{tcolorbox}

% =========================================================================
\section{Steiner / Triple-Point Contributions}
\label{sec:steiner}
% =========================================================================

A \emph{triple point} (also called a Steiner point) arises at a mesh
triangle where three or more cells meet.  The boundary does not pass
cleanly through two VP positions; instead the three boundary segments at
the triple point are required to meet at a shared interior point $\Sv$,
the \emph{Fermat--Torricelli point} of the three VPs
\cite{kuhn1974steiner}.

\subsection{The Fermat--Torricelli point (forward formula)}
\label{sec:steiner-point}

At one triple point let the three VPs be $\vp{1}$, $\vp{2}$, $\vp{3}$
(ordered by the convention stored in \code{pa.tp\_vp\_indices}).

\paragraph{Non-degenerate case (all angles $< 120^\circ$).}
Let
\begin{equation}
  a = \norm{\vp{2}-\vp{3}},\quad
  b = \norm{\vp{1}-\vp{3}},\quad
  c = \norm{\vp{1}-\vp{2}}
\end{equation}
be the side lengths, and let
$A$, $B$, $C$ be the angles of the triangle at $\vp{1}$, $\vp{2}$,
$\vp{3}$ respectively (via the law of cosines).  The barycentric weights
\begin{equation}
  w_i = \frac{a_i}{\sin(A_i + \pi/3)},
  \quad i \in \{1,2,3\},
  \label{eq:steiner-weights}
\end{equation}
(where $a_1 = a$, $a_2 = b$, $a_3 = c$, and $A_i$ is the angle at vertex
$i$) give the Fermat--Torricelli point as the weighted barycentre
\begin{equation}
  \Sv = \frac{w_1\vp{1} + w_2\vp{2} + w_3\vp{3}}{w_1+w_2+w_3}.
  \label{eq:steiner-bary}
\end{equation}

\paragraph{Degenerate case (some angle $\geq 120^\circ$).}
When the angle at vertex $i$ is $\geq 2\pi/3$, the total distance
$\sum_j\norm{\vp{j}-\Sv}$ is minimized by placing $\Sv$ at that vertex:
\begin{equation}
  \Sv = \vp{k}, \qquad k = \arg\max_i A_i
  \quad\text{if } A_k \geq 2\pi/3.
  \label{eq:steiner-degen}
\end{equation}

\begin{tcolorbox}[colback=blue!4!white, colframe=blue!40!black,
                  title={Code: Steiner point}, fontupper=\small]
\codefile{vectorized\_steiner.py}: \code{compute\_steiner\_points(pa)}
\end{tcolorbox}

\subsection{Steiner perimeter value}
\label{sec:steiner-perim-value}

At one triple point with Steiner point $\Sv$, each of the three cells $k$
(indexed by its pair of adjacent VPs $\vp{a}$, $\vp{b}$) contributes a
perimeter value
\begin{equation}
  P_k^{\mathrm{tp}}
    = \norm{\vp{a} - \Sv} + \norm{\vp{b} - \Sv} - \norm{\vp{a} - \vp{b}}.
  \label{eq:steiner-perim-per-cell}
\end{equation}
The total Steiner perimeter is
\begin{equation}
  \Pst(\Lam)
    = \sum_{\text{(tp, cell)}} P_k^{\mathrm{tp}}.
  \label{eq:steiner-perim}
\end{equation}

\begin{tcolorbox}[colback=blue!4!white, colframe=blue!40!black,
                  title={Code: Steiner perimeter value}, fontupper=\small]
\codefile{vectorized\_steiner.py}: \code{compute\_steiner\_perimeter(pa, steiner\_pts)}\\[2pt]
Arrays \code{pa.tp\_contrib\_vp1}, \code{pa.tp\_contrib\_vp2},
\code{pa.tp\_contrib\_tp\_idx} enumerate the (tp, cell) pairs.
\end{tcolorbox}

\subsection{Steiner area value}
\label{sec:steiner-area-value}

Each cell $k$ at a triple point is assigned a \emph{void area} and a
\emph{corner area}:
\begin{align}
  A_{\mathrm{void},k}   &= \mathrm{area}\!\left(\vp{a},\, \vp{b},\, \Sv\right),
  \label{eq:steiner-void-area} \\
  A_{\mathrm{corner},k} &= \mathrm{area}\!\left(\V[v_\mathrm{mesh}],\,
                                                 \vp{a},\, \vp{b}\right),
  \label{eq:steiner-corner-area}
\end{align}
where $\V[v_\mathrm{mesh}]$ is the mesh vertex of the triple-point triangle
that belongs to cell $k$, stored in \code{pa.tp\_contrib\_mesh\_vertex}.
The Steiner area contribution for cell $k$ is
$A_k^{\mathrm{tp}} = A_{\mathrm{void},k} + A_{\mathrm{corner},k}$.

\begin{tcolorbox}[colback=blue!4!white, colframe=blue!40!black,
                  title={Code: Steiner area values}, fontupper=\small]
\codefile{vectorized\_steiner.py}: \code{compute\_steiner\_areas(pa, steiner\_pts)}
\end{tcolorbox}

\subsection{Steiner derivatives}
\label{sec:steiner-fd}

The Steiner perimeter and area contributions are smooth functions of $\Lam$
away from the degenerate $120^\circ$ case.  Their first \emph{and} second
derivatives are computed in closed form, by applying the implicit-function
theorem to the Fermat-point optimality condition.  The full derivation ---
$\partial\Sv/\partial\vp{k}$, $\partial^2\Sv/\partial\vp{k}\partial\vp{l}$,
and the resulting Steiner perimeter and area gradients and Hessians --- is the
subject of its own document,
\codefile{docs/math/03-analytical-steiner-derivatives}.

The earlier finite-difference implementations are retained, renamed with a
\code{\_fd\_reference} suffix, as the independent regression oracle; the
public \code{compute\_steiner\_*} derivative functions dispatch to the
analytical code via the module flag \code{USE\_ANALYTICAL\_STEINER}.

\paragraph{Restricted support.}
Only VPs in the set \code{pa.tp\_affected\_vps} can have non-zero Steiner
derivatives --- the (at most three) VPs at each triple point.

\begin{tcolorbox}[colback=blue!4!white, colframe=blue!40!black,
                  title={Code: Steiner derivatives}, fontupper=\small]
\codefile{vectorized\_steiner.py}: analytical
\code{compute\_steiner\_perimeter\_gradient}, \code{compute\_steiner\_area\_jacobian},
\code{compute\_steiner\_perimeter\_hessian}, \code{compute\_steiner\_area\_hessian},
each with a \code{\_fd\_reference} counterpart.  Full derivation in
\codefile{docs/math/03-analytical-steiner-derivatives}.
\end{tcolorbox}

% =========================================================================
\section{Lagrangian Hessian Assembly}
\label{sec:assembly}
% =========================================================================

\subsection{The Lagrangian}

IPOPT minimizes the augmented objective by maintaining the Lagrangian
\begin{equation}
  \Lagr(\Lam, \mu) = P(\Lam) + \sum_{k=0}^{N-2} \mu_k
                     \bigl(A_k(\Lam) - \Atgt\bigr),
\end{equation}
where $\mu_k$ are Lagrange multipliers supplied by IPOPT at each
iteration.

\subsection{Hessian assembly}

IPOPT's Hessian callback receives the current $\Lam$, the multipliers
$\mu_k$, and an objective scaling factor $\objfac$ (used when the objective
is rescaled internally).  The callback must return
\begin{equation}
  H = \objfac\,\Hess{\Preg} + \objfac\,\Hess{\Pst}
      + \sum_{k=0}^{N-2} \mu_k\,\Hess{A_k^{\mathrm{reg}}}
      + \sum_{k=0}^{N-2} \mu_k\,\Hess{A_k^{\mathrm{st}}}.
  \label{eq:H-assembly}
\end{equation}
Each of the four terms is computed by a separate function and then
summed in the flat sparse-values array:

\begin{center}
\begin{tabular}{lll}
\toprule
Term & Function & Method \\
\midrule
$\Hess{\Preg}$
  & \code{compute\_perimeter\_hessian\_sparse}
  & analytical (§\ref{sec:perimeter-hessian}) \\
$\Hess{\Pst}$
  & \code{compute\_steiner\_perimeter\_hessian}
  & analytical (§\ref{sec:steiner-fd}; doc 03) \\
$\sum_k\mu_k\Hess{A_k^{\mathrm{reg}}}$
  & \code{compute\_area\_hessian\_sparse}
  & analytical (§\ref{sec:area-hessian-1inside}, §\ref{sec:area-2inside}) \\
$\sum_k\mu_k\Hess{A_k^{\mathrm{st}}}$
  & \code{compute\_steiner\_area\_hessian}
  & analytical (§\ref{sec:steiner-fd}; doc 03) \\
\bottomrule
\end{tabular}
\end{center}

\begin{tcolorbox}[colback=blue!4!white, colframe=blue!40!black,
                  title={Code: Hessian callback}, fontupper=\small]
\codefile{perimeter\_optimizer.py}: \code{IPOPTProblemAdapter.hessian()}\\[2pt]
\texttt{return (obj\_factor * (perim\_hess + steiner\_perim\_hess)}\\
\texttt{\phantom{return (}+ area\_hess + steiner\_area\_hess)}
\end{tcolorbox}

\subsection{Lower-triangle convention}
\label{sec:lower-tri}

IPOPT requires only the \emph{lower triangle} of the symmetric matrix $H$,
i.e.\ entries $(i,j)$ with $i \geq j$.  All Hessian functions return a
flat array of values at positions $(\code{pa.hess\_row},\,
\code{pa.hess\_col})$ satisfying \code{hess\_row[k] >= hess\_col[k]}
for all $k$.

\subsection{Sparsity structure}
\label{sec:sparsity}

The non-zero pattern of $H$ (lower triangle) is the union of contributions
from three sources, determined at compile time in
\code{PartitionContour.compile\_arrays()}:

\begin{enumerate}[noitemsep]
  \item \textbf{Perimeter segments.}  Each segment $(a,b)$ contributes the
        three entries $(a,a)$, $(b,b)$, $(\max(a,b),\min(a,b))$.
        Source: \code{pa.seg\_vp1}, \code{pa.seg\_vp2}.

  \item \textbf{Boundary triangles.}  Each triangle with VPs
        $(\mathrm{vp}_1,\mathrm{vp}_2)$ contributes the same three
        entries.  Because the same pair of VPs defines both the perimeter
        segment and the bounding triangle, the perimeter and area patterns
        are identical — no new entries are added.
        Source: \code{pa.btri\_vp1}, \code{pa.btri\_vp2}.

  \item \textbf{Triple points.}  Each triple point has three VPs, giving
        up to six off-diagonal entries (all pairs within the triple point)
        plus three diagonal entries.  Source:
        \code{pa.tp\_vp\_indices}.
\end{enumerate}

% =========================================================================
\section*{Summary table}
% =========================================================================

\begin{center}
\renewcommand{\arraystretch}{1.35}
\begin{tabular}{p{4.2cm}p{4.5cm}p{4.5cm}}
\toprule
\textbf{Quantity} & \textbf{Regular} & \textbf{Steiner / triple-point} \\
\midrule
Perimeter value $P$
  & Analytical \eqref{eq:perimeter}
  & Analytical \eqref{eq:steiner-perim} \\
Perimeter gradient $\grad P$
  & Analytical \eqref{eq:grad-P}
  & Analytical (doc 03) \\
Area Jacobian $J$
  & Analytical \eqref{eq:area-jac-1inside}, \eqref{eq:jac2-1}--\eqref{eq:jac2-2}
  & Analytical (doc 03) \\
Perimeter Hessian $\Hess{P}$
  & Analytical \eqref{eq:H_aa}--\eqref{eq:H_ab}
  & Analytical (doc 03) \\
Area Hessian $\Hess{A_k}$
  & Analytical \eqref{eq:H11-1inside}--\eqref{eq:H12-1inside},
    \eqref{eq:H11-2inside}--\eqref{eq:H22-sub2}
  & Analytical (doc 03) \\
\bottomrule
\end{tabular}
\end{center}

% =========================================================================
\bibliographystyle{plain}
\bibliography{../shared/references}
% =========================================================================

\end{document}
